%% Generated by tools/gen_error_codes.py from docs/errors/ of the compiler;
%% edit the compiler's pages or tools/error_themes.txt, then rerun it.

\clearpage
\section{Exceptions and scope guards}
\label{sec:codes:errors}

The book's chapter \textit{Error handling} specifies the rules behind these codes.

\begin{longtable}{@{}l p{0.78\linewidth} r@{}}
\toprule Code & Message & Page\\ \midrule \endhead
\hyperref[sec:codes:e4003]{E4003} & exception type \errhole{} is already caught by another catching pattern of type \errhole{} & \pageref{sec:codes:e4003}\\
\hyperref[sec:codes:e4012]{E4012} & catching pattern of type \errhole{} catches no exception & \pageref{sec:codes:e4012}\\
\hyperref[sec:codes:e4013]{E4013} & the pattern of a catch scope guard must start with a variable declaration or a field deconstructor & \pageref{sec:codes:e4013}\\
\hyperref[sec:codes:e4021]{E4021} & class type expected, not the class object instance type \errhole{} (i.e. the \& is useless) & \pageref{sec:codes:e4021}\\
\hyperref[sec:codes:e4054]{E4054} & scope guard does not catch exceptions of type \errhole{} & \pageref{sec:codes:e4054}\\
\hyperref[sec:codes:e4055]{E4055} & \errhole{} guard cannot be used when guarding a scope that cannot throw & \pageref{sec:codes:e4055}\\
\hyperref[sec:codes:e4056]{E4056} & \errhole{} scope guard value must be of type void, not \errhole{} & \pageref{sec:codes:e4056}\\
\hyperref[sec:codes:e4145]{E4145} & nothing to catch & \pageref{sec:codes:e4145}\\
\hyperref[sec:codes:e4149]{E4149} & class type \errhole{} does not inherit from exception type \errhole{} & \pageref{sec:codes:e4149}\\
\hyperref[sec:codes:e4206]{E4206} & the function \errhole{} might throw an exception of type \errhole{}, but that is not declared in its prototype & \pageref{sec:codes:e4206}\\
\hyperref[sec:codes:e4207]{E4207} & the method might throw an exception of type \errhole{}, but that is not covered by the method of the ancestor class & \pageref{sec:codes:e4207}\\
\hyperref[sec:codes:e4208]{E4208} & the prototype of the function \errhole{} declares a possible throw of an exception of type \errhole{}, but the function never throws it & \pageref{sec:codes:e4208}\\
\hyperref[sec:codes:e4209]{E4209} & throw statement is not within a function & \pageref{sec:codes:e4209}\\
\hyperref[sec:codes:e4210]{E4210} & cannot throw exception inside a scope guard & \pageref{sec:codes:e4210}\\
\hyperref[sec:codes:e4211]{E4211} & rethrowing an exception of type \errhole{} inside a scope guard is forbidden & \pageref{sec:codes:e4211}\\
\bottomrule
\end{longtable}

\errorcode{E4003}{exception type \errhole{} is already caught by another catching pattern of type \errhole{}}
\label{sec:codes:e4003}

Raised during \textbf{validation}, from \texttt{ValidateErrorMessage::\allowbreak{}ALREADY\_\allowbreak{}CAUGHT} (\textit{src/\allowbreak{}ymirc/\allowbreak{}semantic/\allowbreak{}validator/\allowbreak{}errors.yr}).

\subsubsection*{What it means}

Two catch patterns of the same \tokennolst{try} catch the same exception type, and the second can never run because the first already covers it.

\subsubsection*{How to fix it}

Remove the redundant pattern, or reorder them so the more specific type comes first: a pattern catching a base class shadows every pattern for a derived one written after it.

\errorcode{E4012}{catching pattern of type \errhole{} catches no exception}
\label{sec:codes:e4012}

Raised during \textbf{validation}, from \texttt{ValidateErrorMessage::\allowbreak{}CATCH\_\allowbreak{}NOTHING} (\textit{src/\allowbreak{}ymirc/\allowbreak{}semantic/\allowbreak{}validator/\allowbreak{}errors.yr}).

\subsubsection*{What it means}

A catch pattern names a type no exception thrown inside the guarded block can have, so the pattern can never run.

\subsubsection*{Example}

\textit{test\_\allowbreak{}resources/\allowbreak{}diagnostics/\allowbreak{}test130.yr}:

%% check: skip
\begin{lstlisting}[style=coloredverbatimError]
// a catch pattern for an exception the guarded block cannot throw
class First over Exception {
    pub self () {}
}

class Second over Exception {
    pub self () {}
}

class Other over Exception {
    pub self () {}
}

fn foo (i : i32)-> i32
    throws First, Second;

fn main () {
    let _ = {
        foo (0)
    } catch {
        Other () => { 1 }
        First () => { 2 }
        Second () => { 3 }
    };
}
\end{lstlisting}

\begin{lstlisting}[style=bashVerb, breaklines=true, escapechar=~]
Error[E4012] : catching pattern of type &(test130::Other) catches no exception
 --> test_resources/diagnostics/test130.yr:(21,15)
21  ~\lstbox{\textSFxi{}}~         Other () => { 1 }
    ~\lstbox{\textSFv{}}~               ^
    ~\lstbox{\textSFxi{}}~
    = when validating catching pattern with type &(core::exception::Exception) -- (test_resources/diagnostics/test130.yr:(20,7))
    = when validating test130::main -- (test_resources/diagnostics/test130.yr:(17,4))
    ~\lstbox{\textSFii{}}~~\lstbox{\textSFx{}}~~\lstbox{\textSFx{}}~~\lstbox{\textSFx{}}~~\lstbox{\textSFx{}}~~\lstbox{\textSFx{}}~~\lstbox{\textSFx{}}~~\lstbox{\textSFx{}}~
\end{lstlisting}

\subsubsection*{How to fix it}

Remove the pattern, or widen it to a type the block can actually throw. The types a block may throw are the ones declared by the \tokennolst{throws} of everything it calls.

\errorcode{E4013}{the pattern of a catch scope guard must start with a variable declaration or a field deconstructor}
\label{sec:codes:e4013}

Raised during \textbf{validation}, from \texttt{ValidateErrorMessage::\allowbreak{}CATCH\_\allowbreak{}PATTERN\_\allowbreak{}MULT\_\allowbreak{}OR\_\allowbreak{}VDECL} (\textit{src/\allowbreak{}ymirc/\allowbreak{}semantic/\allowbreak{}validator/\allowbreak{}errors.yr}).

\subsubsection*{What it means}

The pattern of a \tokennolst{catch} does not start with a variable declaration or a field deconstructor. A catch pattern has to bind the exception before it can test anything about it.

\subsubsection*{Example}

\textit{test\_\allowbreak{}resources/\allowbreak{}diagnostics/\allowbreak{}test22.yr}:

%% check: skip
\begin{lstlisting}[style=coloredverbatimError]
fn foo () throws AssertError;

fn main () {
    { foo (); } catch {
        1 => {}
    }
}
\end{lstlisting}

\begin{lstlisting}[style=bashVerb, breaklines=true, escapechar=~]
Error[E4013] : the pattern of a catch scope guard must start with a variable declaration or a field deconstructor
 --> test_resources/diagnostics/test22.yr:(5,9)
 5  ~\lstbox{\textSFxi{}}~         1 => {}
    ~\lstbox{\textSFv{}}~         ^
    ~\lstbox{\textSFxi{}}~
    = when validating catching pattern with type &(core::exception::assertion::AssertError) -- (test_resources/diagnostics/test22.yr:(4,17))
    = when validating test22::main -- (test_resources/diagnostics/test22.yr:(3,4))
    ~\lstbox{\textSFii{}}~~\lstbox{\textSFx{}}~~\lstbox{\textSFx{}}~~\lstbox{\textSFx{}}~~\lstbox{\textSFx{}}~~\lstbox{\textSFx{}}~~\lstbox{\textSFx{}}~~\lstbox{\textSFx{}}~
\end{lstlisting}

\subsubsection*{How to fix it}

Start the pattern with a binding, as in \tokennolst{err: \&SomeException =\textgreater{} \{ ... \}}.

\errorcode{E4021}{class type expected, not the class object instance type \errhole{} (i.e. the \& is useless)}
\label{sec:codes:e4021}

Raised during \textbf{validation}, from \texttt{ValidateErrorMessage::\allowbreak{}CLASS\_\allowbreak{}THROW\_\allowbreak{}PTR} (\textit{src/\allowbreak{}ymirc/\allowbreak{}semantic/\allowbreak{}validator/\allowbreak{}errors.yr}).

\subsubsection*{What it means}

A class \textit{type} was expected and the class \textit{instance} type was given -- the \tokennolst{\&} that turns a class into an instance reference is one level too many here.

\subsubsection*{Example}

\textit{test\_\allowbreak{}resources/\allowbreak{}aka/\allowbreak{}test15.yr}:

%% check: skip
\begin{lstlisting}[style=coloredverbatimError]
class A {
    pub self () {}
}

def B : &A;
\end{lstlisting}

\begin{lstlisting}[style=bashVerb, breaklines=true, escapechar=~]
Error[E4021] : class type expected, not the class object instance type &(test15::A) (i.e. the & is useless)
 --> test_resources/aka/test15.yr:(5,9)
 5  ~\lstbox{\textSFxi{}}~ def B : &A;
    ~\lstbox{\textSFv{}}~         ^
    ~\lstbox{\textSFii{}}~~\lstbox{\textSFx{}}~~\lstbox{\textSFx{}}~~\lstbox{\textSFx{}}~~\lstbox{\textSFx{}}~~\lstbox{\textSFx{}}~~\lstbox{\textSFx{}}~~\lstbox{\textSFx{}}~
\end{lstlisting}

\subsubsection*{How to fix it}

Remove the \tokennolst{\&}.

\errorcode{E4054}{scope guard does not catch exceptions of type \errhole{}}
\label{sec:codes:e4054}

Raised during \textbf{validation}, from \texttt{ValidateErrorMessage::\allowbreak{}EXCEPTION\_\allowbreak{}NOT\_\allowbreak{}CAUGHT} (\textit{src/\allowbreak{}ymirc/\allowbreak{}semantic/\allowbreak{}validator/\allowbreak{}errors.yr}).

\subsubsection*{What it means}

A scope guard is written on a block that can throw an exception the guard does not handle, so the exception would escape past the guard.

\subsubsection*{Example}

\textit{test\_\allowbreak{}resources/\allowbreak{}diagnostics/\allowbreak{}test31.yr}:

%% check: skip
\begin{lstlisting}[style=coloredverbatimError]
class A over Exception { pub self () {} }
class B over Exception { pub self () {} }

fn foo () throws A, B;

fn main () {
    { foo (); } catch {
        A () => {}
    }
}
\end{lstlisting}

\begin{lstlisting}[style=bashVerb, breaklines=true, escapechar=~]
Error[E4054] : scope guard does not catch exceptions of type test31::B
 --> test_resources/diagnostics/test31.yr:(7,11)
 7  ~\lstbox{\textSFxi{}}~     { foo (); } catch {
    ~\lstbox{\textSFv{}}~           ^     ^^^^^
    ~\lstbox{\textSFii{}}~~\lstbox{\textSFx{}}~~\lstbox{\textSFx{}}~~\lstbox{\textSFx{}}~~\lstbox{\textSFx{}}~~\lstbox{\textSFx{}}~~\lstbox{\textSFx{}}~~\lstbox{\textSFx{}}~
\end{lstlisting}

\subsubsection*{How to fix it}

Add a catch for the reported type, or declare the enclosing function as throwing it.

\errorcode{E4055}{\errhole{} guard cannot be used when guarding a scope that cannot throw}
\label{sec:codes:e4055}

Raised during \textbf{validation}, from \texttt{ValidateErrorMessage::\allowbreak{}EXIT\_\allowbreak{}SCOPE\_\allowbreak{}NO\_\allowbreak{}THROW} (\textit{src/\allowbreak{}ymirc/\allowbreak{}semantic/\allowbreak{}validator/\allowbreak{}errors.yr}).

\subsubsection*{What it means}

A \tokennolst{success} or \tokennolst{failure} guard was written on a block that cannot throw. Such a guard distinguishes the two ways of leaving the block, and a block that cannot throw only has one.

\subsubsection*{Example}

\textit{test\_\allowbreak{}resources/\allowbreak{}diagnostics/\allowbreak{}test32.yr}:

%% check: skip
\begin{lstlisting}[style=coloredverbatimError]
fn main () {
    {
        let _ = 1;
    } success {
    }
}
\end{lstlisting}

\begin{lstlisting}[style=bashVerb, breaklines=true, escapechar=~]
Error[E4055] : success guard cannot be used when guarding a scope that cannot throw
 --> test_resources/diagnostics/test32.yr:(4,7)
 4  ~\lstbox{\textSFxi{}}~     } success {
    ~\lstbox{\textSFv{}}~       ^^^^^^^
    ~\lstbox{\textSFii{}}~~\lstbox{\textSFx{}}~~\lstbox{\textSFx{}}~~\lstbox{\textSFx{}}~~\lstbox{\textSFx{}}~~\lstbox{\textSFx{}}~~\lstbox{\textSFx{}}~~\lstbox{\textSFx{}}~
\end{lstlisting}

\subsubsection*{How to fix it}

Use an \tokennolst{exit} guard, which runs however the block is left, or drop the guard. If the block was meant to throw, the call inside it is probably already caught.

\errorcode{E4056}{\errhole{} scope guard value must be of type void, not \errhole{}}
\label{sec:codes:e4056}

Raised during \textbf{validation}, from \texttt{ValidateErrorMessage::\allowbreak{}EXIT\_\allowbreak{}SCOPE\_\allowbreak{}VALUE\_\allowbreak{}TYPE} (\textit{src/\allowbreak{}ymirc/\allowbreak{}semantic/\allowbreak{}validator/\allowbreak{}errors.yr}).

\subsubsection*{What it means}

The body of a scope guard produces a value. A guard runs while a scope is being left, and there is nothing to receive what it would produce.

\subsubsection*{Example}

\textit{test\_\allowbreak{}resources/\allowbreak{}diagnostics/\allowbreak{}test33.yr}:

%% check: skip
\begin{lstlisting}[style=coloredverbatimError]
fn foo () throws AssertError;

fn main () {
    {
        foo ();
    } exit {
        1
    }
}
\end{lstlisting}

\begin{lstlisting}[style=bashVerb, breaklines=true, escapechar=~]
Error[E4056] : exit scope guard value must be of type void, not i32
 --> test_resources/diagnostics/test33.yr:(6,12)
 6  ~\lstbox{\textSFxi{}}~     } exit {
    ~\lstbox{\textSFv{}}~            ^
    ~\lstbox{\textSFii{}}~~\lstbox{\textSFx{}}~~\lstbox{\textSFx{}}~~\lstbox{\textSFx{}}~~\lstbox{\textSFx{}}~~\lstbox{\textSFx{}}~~\lstbox{\textSFx{}}~~\lstbox{\textSFx{}}~
\end{lstlisting}

\subsubsection*{How to fix it}

Make the guard body end in a statement rather than a value -- assign the value to something, or discard it.

\errorcode{E4145}{nothing to catch}
\label{sec:codes:e4145}

Raised during \textbf{validation}, from \texttt{ValidateErrorMessage::\allowbreak{}NOTHING\_\allowbreak{}TO\_\allowbreak{}CATCH} (\textit{src/\allowbreak{}ymirc/\allowbreak{}semantic/\allowbreak{}validator/\allowbreak{}errors.yr}).

\subsubsection*{What it means}

A \tokennolst{catch} was written on a block that cannot throw anything.

\subsubsection*{Example}

\textit{test\_\allowbreak{}resources/\allowbreak{}diagnostics/\allowbreak{}test57.yr}:

%% check: skip
\begin{lstlisting}[style=coloredverbatimError]
fn main () {
    {
        let _ = 1;
    } catch {
        _: &AssertError => {}
    }
}
\end{lstlisting}

\begin{lstlisting}[style=bashVerb, breaklines=true, escapechar=~]
Error[E4145] : nothing to catch
 --> test_resources/diagnostics/test57.yr:(4,7)
 4  ~\lstbox{\textSFxi{}}~     } catch {
    ~\lstbox{\textSFv{}}~       ^^^^^
    ~\lstbox{\textSFii{}}~~\lstbox{\textSFx{}}~~\lstbox{\textSFx{}}~~\lstbox{\textSFx{}}~~\lstbox{\textSFx{}}~~\lstbox{\textSFx{}}~~\lstbox{\textSFx{}}~~\lstbox{\textSFx{}}~
\end{lstlisting}

\subsubsection*{How to fix it}

Remove the catch. If a call inside was expected to throw, check its prototype: it may not declare the exception, or an inner catch may already handle it.

\errorcode{E4149}{class type \errhole{} does not inherit from exception type \errhole{}}
\label{sec:codes:e4149}

Raised during \textbf{validation}, from \texttt{ValidateErrorMessage::\allowbreak{}NOT\_\allowbreak{}AN\_\allowbreak{}EXCEPTION\_\allowbreak{}CLASS} (\textit{src/\allowbreak{}ymirc/\allowbreak{}semantic/\allowbreak{}validator/\allowbreak{}errors.yr}).

\subsubsection*{What it means}

A class is thrown, or caught, that does not derive from the exception base class.

\subsubsection*{Example}

\textit{test\_\allowbreak{}resources/\allowbreak{}diagnostics/\allowbreak{}test58.yr}:

%% check: skip
\begin{lstlisting}[style=coloredverbatimError]
class A { pub self () {} }

fn foo () throws A;

fn main () { foo (); }
\end{lstlisting}

\begin{lstlisting}[style=bashVerb, breaklines=true, escapechar=~]
Error[E4149] : class type test58::A does not inherit from exception type core::exception::Exception
 --> test_resources/diagnostics/test58.yr:(3,11)
 3  ~\lstbox{\textSFxi{}}~ fn foo () throws A;
    ~\lstbox{\textSFv{}}~           ^^^^^^ -
    ~\lstbox{\textSFxi{}}~
    = when validating test58::foo -- (test_resources/diagnostics/test58.yr:(3,4))
    ~\lstbox{\textSFii{}}~~\lstbox{\textSFx{}}~~\lstbox{\textSFx{}}~~\lstbox{\textSFx{}}~~\lstbox{\textSFx{}}~~\lstbox{\textSFx{}}~~\lstbox{\textSFx{}}~~\lstbox{\textSFx{}}~
Error[E4104] : symbol test58::foo has errors
 --> test_resources/diagnostics/test58.yr:(5,14)
 3  ~\lstbox{\textSFxi{}}~ fn foo () throws A;
    ~\lstbox{\textSFv{}}~           ------  error: class type test58::A does not inherit from exception type core::exception::Exception
 4  ~\lstbox{\textSFxi{}}~
 5  ~\lstbox{\textSFxi{}}~ fn main () { foo (); }
    ~\lstbox{\textSFv{}}~              ^^^
    ~\lstbox{\textSFxi{}}~
    = when validating test58::foo -- (test_resources/diagnostics/test58.yr:(3,4))
    = when validating test58::main -- (test_resources/diagnostics/test58.yr:(5,4))
    ~\lstbox{\textSFii{}}~~\lstbox{\textSFx{}}~~\lstbox{\textSFx{}}~~\lstbox{\textSFx{}}~~\lstbox{\textSFx{}}~~\lstbox{\textSFx{}}~~\lstbox{\textSFx{}}~~\lstbox{\textSFx{}}~
\end{lstlisting}

\subsubsection*{How to fix it}

Make the class inherit from the exception type reported in the message.

\errorcode{E4206}{the function \errhole{} might throw an exception of type \errhole{}, but that is not declared in its prototype}
\label{sec:codes:e4206}

Raised during \textbf{validation}, from \texttt{ValidateErrorMessage::\allowbreak{}THROWS\_\allowbreak{}NOT\_\allowbreak{}DECLARED} (\textit{src/\allowbreak{}ymirc/\allowbreak{}semantic/\allowbreak{}validator/\allowbreak{}errors.yr}).

\subsubsection*{What it means}

A function can throw an exception its prototype does not declare. Ymir checks that the throw list of a function covers everything its body can raise.

\subsubsection*{Example}

\textit{test\_\allowbreak{}resources/\allowbreak{}lit\_\allowbreak{}class/\allowbreak{}ctors/\allowbreak{}test7.yr}:

%% check: skip
\begin{lstlisting}[style=coloredverbatimError]
fn foo (t : bool = false)-> i32
    throws AssertError
{
    assert (t);
    12
}

record A {
    let a : i32 = 0;

    pub self ()
        with a = foo ()
    {}

}
\end{lstlisting}

\begin{lstlisting}[style=bashVerb, breaklines=true, escapechar=~]
Error[E4206] : the function test7::A::self might throw an exception of type core::exception::assertion::AssertError, but that is not declared in its prototype
 --> test_resources/lit_class/ctors/test7.yr:(11,9)
11  ~\lstbox{\textSFxi{}}~     pub self ()
    ~\lstbox{\textSFv{}}~         ^^^^
    ~\lstbox{\textSFxi{}}~ Note :
    ~\lstbox{\textSFxi{}}~  --> test_resources/lit_class/ctors/test7.yr:(12,22)
    ~\lstbox{\textSFxi{}}~ 12  ~\lstbox{\textSFxi{}}~         with a = foo ()
    ~\lstbox{\textSFxi{}}~     ~\lstbox{\textSFv{}}~                      ^
    ~\lstbox{\textSFxi{}}~     ~\lstbox{\textSFii{}}~~\lstbox{\textSFx{}}~~\lstbox{\textSFx{}}~~\lstbox{\textSFx{}}~~\lstbox{\textSFx{}}~~\lstbox{\textSFx{}}~~\lstbox{\textSFx{}}~~\lstbox{\textSFx{}}~
    ~\lstbox{\textSFii{}}~~\lstbox{\textSFx{}}~~\lstbox{\textSFx{}}~~\lstbox{\textSFx{}}~~\lstbox{\textSFx{}}~~\lstbox{\textSFx{}}~~\lstbox{\textSFvii{}}~~\lstbox{\textSFx{}}~
\end{lstlisting}

\subsubsection*{How to fix it}

Add the exception to the \tokennolst{throws} clause, or catch it inside the function.

\errorcode{E4207}{the method might throw an exception of type \errhole{}, but that is not covered by the method of the ancestor class}
\label{sec:codes:e4207}

Raised during \textbf{validation}, from \texttt{ValidateErrorMessage::\allowbreak{}THROWS\_\allowbreak{}NOT\_\allowbreak{}DECLARED\_\allowbreak{}OVER} (\textit{src/\allowbreak{}ymirc/\allowbreak{}semantic/\allowbreak{}validator/\allowbreak{}errors.yr}).

\subsubsection*{What it means}

An overriding method throws an exception the ancestor method does not declare, cf. \hyperref[sec:codes:e4186]{E4186}.

\subsubsection*{Example}

\textit{test\_\allowbreak{}resources/\allowbreak{}diagnostics/\allowbreak{}test67.yr}:

%% check: skip
\begin{lstlisting}[style=coloredverbatimError]
class A {
    pub self () {}
    pub fn foo (self) {}
}

class B over A {
    pub self () {}
    pub over foo (self) throws AssertError {
        throw copy AssertError ("x");
    }
}

fn main () { let _ = copy B (); }
\end{lstlisting}

\begin{lstlisting}[style=bashVerb, breaklines=true, escapechar=~]
Error[E4186] : the throwers of the overriding method test67::B::foo ()-> void are not compatible with the throwers of the ancestor method test67::A::foo ()-> void
 --> test_resources/diagnostics/test67.yr:(3,12)
 3  ~\lstbox{\textSFxi{}}~     pub fn foo (self) {}
    ~\lstbox{\textSFv{}}~            ^^^
 --> test_resources/diagnostics/test67.yr:(8,32)
 8  ~\lstbox{\textSFxi{}}~     pub over foo (self) throws AssertError {
    ~\lstbox{\textSFv{}}~                                -----------  error: the method might throw an exception of type core::exception::assertion::AssertError, but that is not covered by the method of the ancestor class
    ~\lstbox{\textSFxi{}}~
    = when validating -- (test_resources/diagnostics/test67.yr:(6,7))
    ~\lstbox{\textSFii{}}~~\lstbox{\textSFx{}}~~\lstbox{\textSFx{}}~~\lstbox{\textSFx{}}~~\lstbox{\textSFx{}}~~\lstbox{\textSFx{}}~~\lstbox{\textSFx{}}~~\lstbox{\textSFx{}}~
Error[E4104] : symbol test67::B has errors
 --> test_resources/diagnostics/test67.yr:(13,27)
13  ~\lstbox{\textSFxi{}}~ fn main () { let _ = copy B (); }
    ~\lstbox{\textSFv{}}~                           ^
 --> test_resources/diagnostics/test67.yr:(3,12)
 3  ~\lstbox{\textSFxi{}}~     pub fn foo (self) {}
    ~\lstbox{\textSFv{}}~            ---  error: the throwers of the overriding method test67::B::foo ()-> void are not compatible with the throwers of the ancestor method test67::A::foo ()-> void
   ...
 8  ~\lstbox{\textSFxi{}}~     pub over foo (self) throws AssertError {
    ~\lstbox{\textSFv{}}~                                -----------  error: the method might throw an exception of type core::exception::assertion::AssertError, but that is not covered by the method of the ancestor class
    ~\lstbox{\textSFxi{}}~
    = when validating -- (test_resources/diagnostics/test67.yr:(6,7))
    = when validating test67::main -- (test_resources/diagnostics/test67.yr:(13,4))
    ~\lstbox{\textSFii{}}~~\lstbox{\textSFx{}}~~\lstbox{\textSFx{}}~~\lstbox{\textSFx{}}~~\lstbox{\textSFx{}}~~\lstbox{\textSFx{}}~~\lstbox{\textSFx{}}~~\lstbox{\textSFx{}}~
\end{lstlisting}

\subsubsection*{How to fix it}

Catch it in the overriding method, or declare it on the ancestor.

\errorcode{E4208}{the prototype of the function \errhole{} declares a possible throw of an exception of type \errhole{}, but the function never throws it}
\label{sec:codes:e4208}

Raised during \textbf{validation}, from \texttt{ValidateErrorMessage::\allowbreak{}THROWS\_\allowbreak{}NOT\_\allowbreak{}USED} (\textit{src/\allowbreak{}ymirc/\allowbreak{}semantic/\allowbreak{}validator/\allowbreak{}errors.yr}).

\subsubsection*{What it means}

A function declares that it throws an exception its body can never raise. The declaration forces every caller to handle something that cannot happen.

\subsubsection*{Example}

\textit{test\_\allowbreak{}resources/\allowbreak{}lit\_\allowbreak{}class/\allowbreak{}ctors/\allowbreak{}test8.yr}:

%% check: skip
\begin{lstlisting}[style=coloredverbatimError]
fn foo (t : bool = false)-> i32 {
    t;
    12
}

record A {
    let a : i32 = 0;

    pub self ()
        with a = foo ()
        throws AssertError
    {}

}
\end{lstlisting}

\begin{lstlisting}[style=bashVerb, breaklines=true, escapechar=~]
Error[E4208] : the prototype of the function test8::A::self declares a possible throw of an exception of type core::exception::assertion::AssertError, but the function never throws it
 --> test_resources/lit_class/ctors/test8.yr:(9,9)
 9  ~\lstbox{\textSFxi{}}~     pub self ()
    ~\lstbox{\textSFv{}}~         ^^^^
    ~\lstbox{\textSFxi{}}~ Note :
    ~\lstbox{\textSFxi{}}~  --> test_resources/lit_class/ctors/test8.yr:(11,16)
    ~\lstbox{\textSFxi{}}~ 11  ~\lstbox{\textSFxi{}}~         throws AssertError
    ~\lstbox{\textSFxi{}}~     ~\lstbox{\textSFv{}}~                ^^^^^^^^^^^
    ~\lstbox{\textSFxi{}}~     ~\lstbox{\textSFii{}}~~\lstbox{\textSFx{}}~~\lstbox{\textSFx{}}~~\lstbox{\textSFx{}}~~\lstbox{\textSFx{}}~~\lstbox{\textSFx{}}~~\lstbox{\textSFx{}}~~\lstbox{\textSFx{}}~
    ~\lstbox{\textSFii{}}~~\lstbox{\textSFx{}}~~\lstbox{\textSFx{}}~~\lstbox{\textSFx{}}~~\lstbox{\textSFx{}}~~\lstbox{\textSFx{}}~~\lstbox{\textSFvii{}}~~\lstbox{\textSFx{}}~
\end{lstlisting}

\subsubsection*{How to fix it}

Remove the exception from the \tokennolst{throws} clause.

\errorcode{E4209}{throw statement is not within a function}
\label{sec:codes:e4209}

Raised during \textbf{validation}, from \texttt{ValidateErrorMessage::\allowbreak{}THROW\_\allowbreak{}NO\_\allowbreak{}FUNCTION} (\textit{src/\allowbreak{}ymirc/\allowbreak{}semantic/\allowbreak{}validator/\allowbreak{}errors.yr}).

\subsubsection*{What it means}

A \tokennolst{throw} was written outside any function.

\subsubsection*{Example}

\textit{test\_\allowbreak{}resources/\allowbreak{}diagnostics/\allowbreak{}test70.yr}:

%% check: skip
\begin{lstlisting}[style=coloredverbatimError]
lazy X : i32 = { throw copy AssertError ("x"); 1 };

fn main () { X; }
\end{lstlisting}

\begin{lstlisting}[style=bashVerb, breaklines=true, escapechar=~]
Error[E4209] : throw statement is not within a function
 --> test_resources/diagnostics/test70.yr:(1,18)
 1  ~\lstbox{\textSFxi{}}~ lazy X : i32 = { throw copy AssertError ("x"); 1 };
    ~\lstbox{\textSFv{}}~                  ^^^^^
    ~\lstbox{\textSFxi{}}~
    = when validating -- (test_resources/diagnostics/test70.yr:(1,6))
    ~\lstbox{\textSFii{}}~~\lstbox{\textSFx{}}~~\lstbox{\textSFx{}}~~\lstbox{\textSFx{}}~~\lstbox{\textSFx{}}~~\lstbox{\textSFx{}}~~\lstbox{\textSFx{}}~~\lstbox{\textSFx{}}~
Error[E4104] : symbol test70::X has errors
 --> test_resources/diagnostics/test70.yr:(3,14)
 1  ~\lstbox{\textSFxi{}}~ lazy X : i32 = { throw copy AssertError ("x"); 1 };
    ~\lstbox{\textSFv{}}~                  -----  error: throw statement is not within a function
 2  ~\lstbox{\textSFxi{}}~
 3  ~\lstbox{\textSFxi{}}~ fn main () { X; }
    ~\lstbox{\textSFv{}}~              ^
    ~\lstbox{\textSFxi{}}~
    = when validating -- (test_resources/diagnostics/test70.yr:(1,6))
    = when validating test70::main -- (test_resources/diagnostics/test70.yr:(3,4))
    ~\lstbox{\textSFii{}}~~\lstbox{\textSFx{}}~~\lstbox{\textSFx{}}~~\lstbox{\textSFx{}}~~\lstbox{\textSFx{}}~~\lstbox{\textSFx{}}~~\lstbox{\textSFx{}}~~\lstbox{\textSFx{}}~
\end{lstlisting}

\subsubsection*{How to fix it}

Move it into the function it belongs to.

\errorcode{E4210}{cannot throw exception inside a scope guard}
\label{sec:codes:e4210}

Raised during \textbf{validation}, from \texttt{ValidateErrorMessage::\allowbreak{}THROW\_\allowbreak{}SCOPE\_\allowbreak{}GUARD} (\textit{src/\allowbreak{}ymirc/\allowbreak{}semantic/\allowbreak{}validator/\allowbreak{}errors.yr}).

\subsubsection*{What it means}

A \tokennolst{throw} inside a scope guard would escape the guard uncaught. A guard runs while a scope is being left, so an exception leaving it would interrupt an unwinding already in progress.

\subsubsection*{Example}

\textit{test\_\allowbreak{}resources/\allowbreak{}scope\_\allowbreak{}guards/\allowbreak{}test10.yr}:

%% check: skip
\begin{lstlisting}[style=coloredverbatimError]
fn main () {
    {

    } exit {
        throw copy AssertError ("");
    }
}
\end{lstlisting}

\begin{lstlisting}[style=bashVerb, breaklines=true, escapechar=~]
Error[E4210] : cannot throw exception inside a scope guard
 --> test_resources/scope_guards/test10.yr:(4,7)
 4  ~\lstbox{\textSFxi{}}~     } exit {
    ~\lstbox{\textSFv{}}~       ^^^^
   ...
 --> test_resources/scope_guards/test10.yr:(5,9)
 5  ~\lstbox{\textSFxi{}}~         throw copy AssertError ("");
    ~\lstbox{\textSFv{}}~         ^^^^^
    ~\lstbox{\textSFii{}}~~\lstbox{\textSFx{}}~~\lstbox{\textSFx{}}~~\lstbox{\textSFx{}}~~\lstbox{\textSFx{}}~~\lstbox{\textSFx{}}~~\lstbox{\textSFx{}}~~\lstbox{\textSFx{}}~
\end{lstlisting}

\subsubsection*{How to fix it}

Catch the exception inside the guard. A \tokennolst{throw} in a nested block whose own \tokennolst{catch} handles every thrown type does not cross the guard and is accepted.

\errorcode{E4211}{rethrowing an exception of type \errhole{} inside a scope guard is forbidden}
\label{sec:codes:e4211}

Raised during \textbf{validation}, from \texttt{ValidateErrorMessage::\allowbreak{}THROW\_\allowbreak{}SCOPE\_\allowbreak{}GUARD\_\allowbreak{}RETHROW} (\textit{src/\allowbreak{}ymirc/\allowbreak{}semantic/\allowbreak{}validator/\allowbreak{}errors.yr}).

\subsubsection*{What it means}

A caught exception is rethrown from inside a scope guard, which would let it escape the guard, cf. \hyperref[sec:codes:e4210]{E4210}.

\subsubsection*{Example}

\textit{test\_\allowbreak{}resources/\allowbreak{}scope\_\allowbreak{}guards/\allowbreak{}test11.yr}:

%% check: skip
\begin{lstlisting}[style=coloredverbatimError]
fn foo ()
    throws AssertError;

fn main () {
    {

    } exit {
        foo ();
    }
}
\end{lstlisting}

\begin{lstlisting}[style=bashVerb, breaklines=true, escapechar=~]
Error[E4211] : rethrowing an exception of type core::exception::assertion::AssertError inside a scope guard is forbidden
 --> test_resources/scope_guards/test11.yr:(7,7)
 7  ~\lstbox{\textSFxi{}}~     } exit {
    ~\lstbox{\textSFv{}}~       ^^^^
   ...
 --> test_resources/scope_guards/test11.yr:(8,13)
 8  ~\lstbox{\textSFxi{}}~         foo ();
    ~\lstbox{\textSFv{}}~             ^
    ~\lstbox{\textSFii{}}~~\lstbox{\textSFx{}}~~\lstbox{\textSFx{}}~~\lstbox{\textSFx{}}~~\lstbox{\textSFx{}}~~\lstbox{\textSFx{}}~~\lstbox{\textSFx{}}~~\lstbox{\textSFx{}}~
\end{lstlisting}

\subsubsection*{How to fix it}

Handle the exception inside the guard rather than rethrowing it.
